%! The Normal Distribution and z-Scores — Unit 1, Lesson 1.7 — Lesson Plan
% GRADUAL-RELEASE plan (warm-up / guided notes and practice / debrief / close and start homework).
% REGENERATED from the retired EFFL plan: there is no activity, no spoiler rule, and no
% QuickNotes. The guided notes carry the release row by row in one Main Ideas / Notes table:
% I do (each row's definition lines and first problem) -> we do (the rest of the grid, then the
% Guided Practice row). THERE IS NO INDEPENDENT PRACTICE SET — the homework is the individual
% practice, and the period ends with students starting it. Teacher-facing. Every box is
% `skillbox` (breakable) with a \boxguard ahead of it; this course has no `fixedskillbox`,
% no tier boxes, and no exit ticket.
% This is the LAST content lesson of Unit 1 — the Close previews the unit test and Unit 2.
\documentclass[10pt]{article}
\usepackage{apstats-article}
\usepackage{apstats-boxes}

\newcommand{\UnitNumberName}{Unit 1: Exploring One-Variable Data \quad}
\newcommand{\LessonNumberName}{Lesson 1.7: The Normal Distribution and z-Scores}

\begin{document}
\begin{center}
    {\Large\bfseries \CourseName \\
    {\small\bfseries \UnitNumberName \LessonNumberName}}
\end{center}

\begin{tcolorbox}[colback=sky, colframe=navy, boxrule=0.9pt, arc=2mm,
    left=2mm, right=2mm, top=1.5mm, bottom=1.5mm]
    \textbf{Primary Objective:} Students will distinguish a parameter from a statistic, compute
    and interpret a $z$-score in context, use the empirical rule to find approximate proportions
    under a normal model, and compare the relative positions of values drawn from two different
    distributions.
    \\[2pt]
    \textbf{CED alignment:} Topics 1.9, 1.10 \quad$\cdot$\quad Big Idea: VAR (Variation and
    Distribution) \quad$\cdot$\quad Skills 2.D, 3.A \quad$\cdot$\quad LO VAR-2.A, VAR-2.B,
    VAR-2.C \quad$\cdot$\quad EK VAR-2.A.1--2.A.4, 2.B.1--2.B.4, 2.C.1
    \\[2pt]
    \textbf{Lesson model:} \emph{Gradual release --- I do, we do, you do --- all of it inside
    the guided notes.} There is no group activity. The notes name the vocabulary as they teach
    it and deliver the instruction row by row in one \emph{Main Ideas / Notes} table --- each
    idea a definition to complete and a grid of short problems, the first worked by you and the
    rest by the class; the last row, Guided Practice, is worked together with students holding
    the pen; \textbf{the homework is the individual practice}, and students start it in class
    while you circulate. The whole-class debrief is \emph{spoken} --- there is no transfer set
    and no reflection box on the page.
    \\[2pt]
    \textbf{Note on the ``no sketch'' rule:} every dotplot, histogram, and normal curve in this
    lesson is \textbf{given}. Students read, annotate, and compute from them. Do not add a ``draw
    the normal curve'' step --- the exam asks students to \emph{use} the model, and hand-drawing
    a bell shape teaches nothing about $\mu$, $\sigma$, or the empirical rule.
\end{tcolorbox}

\renewcommand{\arraystretch}{1.4}
\boxguard
\begin{skillbox}[Learning Targets \& Key Understandings]{goldbox}
\begin{tabularx}{\linewidth}{@{}X|X@{}}
\textbf{Learning targets (student language)} & \textbf{Key understandings (the ``why'')} \\[2pt]
\hline
\vspace{-6pt}
\begin{itemize}[leftmargin=1.1em, nosep, after=\vspace{2pt}]
  \item I can tell a \textbf{parameter} from a \textbf{statistic} and write the right symbol for
        each.
  \item I can compute a \textbf{$z$-score} and interpret it in context, sign included.
  \item I can use the \textbf{empirical rule} on an approximately \textbf{normal} distribution,
        including halving a tail.
  \item I can use $z$-scores to compare two values that came from two different distributions.
\end{itemize}
&
\vspace{-6pt}
\begin{itemize}[leftmargin=1.1em, nosep, after=\vspace{2pt}]
  \item A \textbf{parameter} summarizes a population ($\mu$, $\sigma$); a \textbf{statistic}
        summarizes a sample ($\bar x$, $s_x$). Same idea, different object (EK VAR-2.A.1).
  \item A normal curve is mound-shaped and symmetric and is fixed entirely by $\mu$ and $\sigma$
        (EK VAR-2.A.2, 2.A.4).
  \item $z = (x_i - \mu)/\sigma$ counts standard deviations, which is a \emph{unitless} count ---
        that is exactly why it can compare grams to seconds (EK VAR-2.B.1, 2.B.2, 2.C.1).
  \item \textbf{Target misconception:} that the bigger raw number is the better result --- that a
        raw score can be compared across two different distributions without standardizing.
        \textbf{Notes problem~11 is the trap} and \textbf{problem~13 its companion}: a negative
        $z$ read as automatically worse, when in a race a lower time is better.
  \item \textbf{Second misconception:} that 68--95--99.7 applies to any data set. It applies to a
        \emph{mound-shaped, symmetric} one; a skewed or two-peaked display voids it
        (EK VAR-2.A.3) --- notes problem~14.
\end{itemize}
\end{tabularx}
\end{skillbox}

\boxguard
\begin{skillbox}[{Vocabulary, Concepts, \& Theorems --- taught directly in the guided notes}]{greenbox}
\begin{minipage}[t]{0.48\linewidth}
    \begin{itemize}
        \item \textbf{Parameter} --- a numerical summary of a \emph{population} (EK VAR-2.A.1)
        \item \textbf{$\mu$, $\sigma$} --- the population mean and population standard
              deviation; contrast with $\bar x$ and $s_x$ from Lesson 1.5
        \item \textbf{Normal distribution / normal curve} --- mound-shaped and symmetric, fixed
              by $\mu$ and $\sigma$ (EK VAR-2.A.2)
        \item \textbf{Approximately normal} --- the honest phrase; real data is never exactly
              normal (EK VAR-2.A.4)
    \end{itemize}
\end{minipage}
\hfill
\begin{minipage}[t]{0.48\linewidth}
    \begin{itemize}
        \item \textbf{Empirical rule (68--95--99.7)} --- approximate proportions within 1, 2, and
              3 standard deviations of $\mu$ (EK VAR-2.A.3)
        \item \textbf{Standardized score / $z$-score} --- $z = (x_i - \mu)/\sigma$
              (EK VAR-2.B.1, 2.B.2)
        \item \textbf{Relative position} --- $z$-scores and percentiles compare points within or
              between data sets (EK VAR-2.C.1)
        \item \textbf{Carried forward} --- mean and standard deviation (1.5); percentile and the
              five-number summary (1.6); shape, centre, spread (1.4)
    \end{itemize}
\end{minipage}
\end{skillbox}

% A tabularx never splits, so guard generously ahead of it.
\boxguard[30]
\begin{skillbox}[Lesson at a Glance --- \MeetingLength]{sky}
\renewcommand{\arraystretch}{1.25}
\begin{tabularx}{\linewidth}{@{}l c X X@{}}
\textbf{Phase} & \textbf{Min} & \textbf{Students} & \textbf{Teacher} \\
\hline
Warm-Up & 5 & Two five-number summaries, counting steps away from a mean, and a percent split between two ends & Debrief item~3 aloud; leave the ``split it evenly'' line on the board all period \\
Guided Notes \& Practice & 34 & Fill the vocabulary box and the notes table: problems 1--14 row by row, then \emph{The Bread Loaf} (15--20) together, pen in their hand & \textbf{I do} each row's prose, work blocks, and first problem, \textbf{we do} the rest of its grid (24) $\to$ \textbf{we do} Guided Practice (10) \\
Debrief & 8 & Pens down, papers up --- answer aloud, nothing to fill in & Open on an \emph{almost}-right answer to problem~12, then run problems 11, 13, and~14 public \\
Close \& Assign & 8 & \textbf{Start the homework in class}, alone and silent & Announce the paper homework and the unit-test date, circulate for your formative read \\
\end{tabularx}
\end{skillbox}

\boxguard[20]
\begin{skillbox}[Warm-Up --- Activate Prior Knowledge (5 min)]{sky}
\begin{minipage}[t]{0.48\linewidth}
    \textbf{The three items and what each seeds}
    \begin{itemize}
        \item \textbf{1.} Two five-number summaries, two IQRs. \textbf{Spirals} Lesson 1.6 ---
              and quietly re-establishes that ``more spread out'' is a measurable thing.
        \item \textbf{2.} A mean of 20 and a standard deviation of 3; how many 3-point steps
              from the mean is 26? is 17? \textbf{Seeds:} the $z$-score arithmetic of notes
              row~2, done once with numbers that divide evenly.
        \item \textbf{3.} 360 of 400 inside a range; what percent is outside, and what percent
              past the high end? \textbf{Seeds:} the halve-the-tail move that notes problems
              16--17 rest on.
    \end{itemize}
\end{minipage}
\hfill
\begin{minipage}[t]{0.48\linewidth}
    \textbf{Running it}
    \begin{itemize}
        \item Item~2 uses numbers that divide evenly on purpose --- the arithmetic must not be
              the obstacle when the notes name it as $z$ eight minutes later. If a student asks
              what ``2 steps'' \emph{means}, say ``that is the first thing we name today.''
        \item \textbf{Debrief item~3 aloud} and leave \emph{``symmetric, so split the leftovers
              evenly''} on the board for the whole period, in the students' own words. Notes
              problems 16, 17, and the AP practice all reach back for it.
        \item The percentages in item~3 are 90 / 10 / 5 rather than 95 / 5 / 2.5 --- keep it
              that way, so that the empirical rule's real numbers arrive in notes problem~7
              out of the students' own band counts.
        \item Four minutes, five at the outside. Do not let item~1 reopen the boxplot
              conversation.
    \end{itemize}
\end{minipage}
\end{skillbox}

\boxguard
\begin{skillbox}[Guided Notes \& Practice --- I do, then we do (34 min)]{sky}
\textbf{This is the whole lesson.} There is no group activity, and there is no independent
practice set on this page: \textbf{the homework is the individual practice}, started in class.
The notes are one \emph{Main Ideas / Notes} table --- five rows, 20 short problems. \textbf{In
every row you narrate the prose, fill the lines, work the block, and work the first problem of the
grid; the class works the rest, pen in their hand.} The vocabulary box is filled \emph{as you go},
never in advance.
\begin{multicols}{2}
\textbf{Parameter \& Statistic (5 min).} All ten fills, then the four symbols on the board in
two columns --- $\mu$, $\sigma$ beside $\bar x$, $s_x$ --- and leave them up for the rest of the
period; this contrast pays off in every remaining unit. Work problem~1 yourself; 2 and~3 are
theirs. \textbf{Problem~3 is the one to check}: every one of the 400 temperatures was recorded,
so the symbol is $\sigma$.

\textbf{$z$-score (9 min).} The definition line, the formula, then the three steps. The two
dotplots and the field summary table are already on the page --- point at Maya's gold dot and
Dev's, and say out loud that 12.0 metres and 12.3 seconds are not on one scale. \textbf{You work
both $z$-score blocks at the board} so the arithmetic is settled in public; then 4--6 are theirs.
\textbf{Problem~5 is where the sign becomes meaning}: insist on the direction \emph{and} on
``good in a sprint.'' Problem~6 is $z = 0$.

\textbf{Normal distribution \& Empirical rule (10 min).} Fill the shape lines and the three
percentages, pointing at the given curve. Work the loaves' band block yourself --- 341/476/499
out of 500 --- then \textbf{problem~7 is the arithmetic that earns the rule}: they do the clinic's
400 readings the same way, 273/381/399. \textbf{Do not announce 68--95--99.7 before that column is
filled.} Their own numbers are 68.2, 95.2, 99.8 and 68.25, 95.25, 99.75; problem~9 rounds them
into the rule. Problem~8 is the point: shape, not units, is what fixes the percentages.

\columnbreak
\textbf{Whose day was better? (5 min) --- the crux.} Read the coach's argument aloud exactly as
written, then fill the boxed sentence. \textbf{Problem~11 is the trap} --- take a hand count on
``is the coach right?'' before anyone speaks. Work nothing here yourself; this is the row where
you find out who is with you. \textbf{Problem~12 settles it with the two $z$-scores} and must end
with a named athlete and a reason. \textbf{Problem~13 is the companion error} (negative read as
worse) and \textbf{problem~14 is the counterweight} --- a skewed or two-peaked set does not obey
the rule.

\textbf{We do (about 10 min): The Bread Loaf.} Problems 15--20, worked entirely together, on the
same bakery the notes already pictured --- so no new context has to be loaded in ten minutes. The
three work blocks (834 g, the halved tail, 754 g) are worked with the class, not for them.
Students hold the pen; you hold the questions: \emph{``What are the units on 1.7?''} \quad
\emph{``840 is how many standard deviations above 800? So what does the rule already tell you?''}
\quad \emph{``Is $-2.3$ unusual, or just below average? Which word does the evidence support?''}
\textbf{Problem~19 is the judgment problem} and \textbf{20 is the honest limit} --- protect time
for both; if you only get 15, 16, 19, and~20, that is the right trade.

Circulate and demand the reason, not the number:
\begin{itemize}[leftmargin=1.1em, nosep]
  \item \textbf{Problem~11, the crux:} ``Write the unit after each number and read your sentence
        back to me.''
  \item \textbf{Problem~13:} ``Below the mean \emph{time}. Faster or slower than the field?''
  \item \textbf{Problem~17:} ``Five percent is outside on \emph{both} ends. How much on the end
        you were asked about?''
\end{itemize}
While circulating, pick the papers for the debrief --- one \emph{almost}-right problem~12 (both
$z$-scores computed, direction dropped) and one problem~14. \textbf{Your formative read comes
twice}: here, and again in the last eight minutes as students start the homework.
\end{multicols}
\end{skillbox}

\boxguard
\begin{skillbox}[Debrief --- whole class (8 min)]{sky}
Pens down, papers up. \textbf{This debrief is spoken} --- there is no box at the end of the
notes to fill in, so it lives entirely on the board and in the room.
\begin{multicols}{2}
\begin{enumerate}[leftmargin=1.3em, nosep, itemsep=3pt]
  \item \textbf{Open on the \emph{almost}-right problem~12} you collected: both $z$-scores
        right, the direction dropped. Put it up and ask the room what is missing. That one
        omission is the difference between a scored answer and an unscored one on the exam.
  \item Then problem~11, the trap, in public: metres and seconds, two fields, never one scale.
        Write $+3.0$ m and $-0.9$ s beside $2.0$ and $-1.5$ and make someone say why only the
        second pair can be compared.
  \item Problem~13 next --- the sign. \emph{``Would you rather be above or below average in a
        100-metre time?''} Take a hand count, then settle it.
  \item Problem~14 last, off their own problem~7 numbers: the rule is a fact \emph{about a
        shape}. Ask what the bread and the thermometers had in common, and what a skewed record
        would do to those three percentages.
\end{enumerate}
\columnbreak
\textbf{Two things to demand aloud.} First, \textbf{make someone state a $z$-score interpretation
in full} --- value, how many standard deviations, which side, and the variable named. ``$z = -1.5$''
earns nothing on its own. Second, \textbf{make them generate a variable where a negative
$z$-score is good news} and one where it is bad; that is where transfer shows up now that there
is no transfer set on paper.

\textbf{One closing line only} (EK VAR-2.B.3, 2.B.4): technology --- a calculator or a standard
normal table --- turns a $z$-score into an exact proportion. \textbf{We do not do that today.} It
returns in Unit~5.

\textbf{The sentence to land:} \emph{Points, seconds, grams, and degrees cannot be compared.
Standard deviations can. Once you count in standard deviations, every distribution speaks the
same language --- and if it is mound-shaped, it always says 68, 95, and 99.7.} Say it, write it,
and point at Maya's 2.0 and Dev's $-1.5$ while you do.

\textbf{If the debrief runs short}, cut item~4 to one sentence --- item~1 is the one that cannot
be cut. \textbf{Do not let the debrief eat the last eight minutes.}
\end{multicols}
\end{skillbox}

\boxguard
\begin{skillbox}[AP Practice --- assigned, not scored]{goldbox}
Four AP-style multiple-choice items on newborn birth weights ($\mu = 3{,}300$ g,
$\sigma = 500$ g, approximately normal, with the curve given and its band boundaries marked) ---
a straight $z$-score computation, an empirical-rule proportion that requires \emph{halving} a
tail, a comparison across two different distributions (the lesson's crux in AP clothing), and
\textbf{one item whose correct answer is that the empirical rule does not apply}, because the
given length-of-stay histogram is strongly right-skewed. That last one is a real AP item type and
the honest limit of the model. Then one multi-part free response on the same birth-weight record:
compute and interpret a $z$-score, use the empirical rule, compare relative position against a
second hospital's distribution, and \textbf{part (d) spirals to Lessons 1.4 and 1.6} --- judge
from a display whether a normal model is reasonable at all, naming the shape features that decide
it.

\textbf{Not scored --- the cover's score column reads \textbf{NA} for this page.} Go over it in
the unit-test review and hold the AP standard out loud: \emph{in context or it does not count}.
``$z = -1.5$'' earns nothing on its own; ``this newborn weighs 1.5 standard deviations less than
the mean birth weight of 3,300 g'' is the answer.
\end{skillbox}

\boxguard
\begin{skillbox}[Homework --- scored, due the first class after two study halls]{goldbox}
Five exercises in a third context --- heart rate, two tests, and a bus route: \textbf{1}
parameter or statistic, on four given summaries; \textbf{2} three $z$-scores from one
distribution, including $z = 0$ and a negative one interpreted in a sentence; \textbf{3} the crux
head-on --- two students, two tests, different means and standard deviations, higher raw score
loses; \textbf{4} the empirical rule off a given normal curve, including a tail proportion;
\textbf{5} the AP-style multiple-choice item with a one-sentence justification. The packet closes
with the unit-test and Unit~2 preview.

\textbf{Start it in the last eight minutes of the period, in class and alone.} \textbf{Scoring:}
10 points --- 2 per item. Eight minutes will not finish five items, and that is by design: point
them at 1, 2, and~3, where a wrong start is cheap to catch; 4 and~5 travel home cleanly.
\textbf{Item~3 is the one that predicts the unit test} --- a student who answers ``Ben, he scored
88'' has not made the turn this lesson exists to make.

\textbf{Paper homework is the default for this course} --- assign this page unless you have
decided otherwise. \textbf{DeltaMath override (occasional):} on the days you assign DeltaMath
instead, it replaces this page, you strike through the score cell on the cover, and this page
stays in the packet as extra practice. Never assign both.
\end{skillbox}

\boxguard
\begin{skillbox}[Watch For (while circulating)]{redbox}
\begin{itemize}[leftmargin=1.2em, nosep]
  \item \textbf{``Dev, because 12.3 is bigger''} (notes~11, HW~3 --- the target misconception).
        Do not correct it; make it collide with the units. Probe: \emph{``Bigger what? Write the
        unit after each number and read it back to me.''}
  \item \textbf{Stopping at the difference} (notes 4--6, HW~3): $+3.0$ versus $-0.9$ feels
        decisive. Probe: \emph{``Is 3.0 metres a lot in the shot put? Look at the dotplot before
        you answer.''}
  \item \textbf{Sign thrown away} (notes 4--6, HW~2, HW~5, AP~3): students report Dev as
        ``1.5'' and lose the whole sprint insight. Probe: \emph{``Faster or slower than average?
        So which side of the mean?''}
  \item \textbf{Lower-is-better inverted} (notes~13, HW~5): having learned that negative means
        below, students conclude Dev did badly. Probe: \emph{``Would you rather be above or below
        average in a 100-metre time?''}
  \item \textbf{68--95--99.7 applied to anything} (notes~14, AP~4): the second misconception, and
        the one the exam tests. Probe: \emph{``What did the bread and the thermometers have in
        common? Does this one have it?''}
  \item \textbf{The whole tail instead of half} (notes 16--17, HW~4, AP~2): students answer 5\%
        where 2.5\% is wanted. Probe: \emph{``Five percent is outside on \emph{both} ends. How
        much is on the end you were asked about?''}
  \item \textbf{$\mu$ and $\bar x$ used interchangeably} (notes 1--3, HW~1). Probe:
        \emph{``Did somebody measure every loaf, or a handful of loaves?''}
\end{itemize}
Cold-call prompts: \emph{``Give me a value with a $z$-score of zero.''} \quad \emph{``Name a
variable where a negative $z$-score is good news.''}
\end{skillbox}

\boxguard
\begin{skillbox}[Close \& Start the Homework (8 min)]{goldbox}
\textbf{Students start the homework here, in class, alone and silent --- this is the individual
practice, and the first minutes of it are your best formative read.} Hand the launch first ---
five exercises on paper, scored, \textbf{due the first class after two study halls} --- and say
plainly that the AP Practice page before it is theirs to work and bring to the unit-test review.
\textbf{Then name the unit-test date}, because this is the last content lesson of Unit~1.
Circulate: item~3 is the column that tells you whether the crux row landed. Sort what you see
into three piles as you walk --- correct with both $z$-scores (ready); correct answer justified
by the raw scores (ready, needs the language); ``Ben, he scored 88'' (the misconception survived)
--- and \textbf{reopen Maya and Dev in the unit-test review} if that third pile ran past a third
of the room. Close on what changed: \emph{``You walked in believing the bigger number wins. You
walk out knowing that a number means nothing until you know what it is being compared to --- and
that once you count in standard deviations, grams and seconds finally sit on the same page.''}
\textbf{Preview:} the \emph{Unit~1 test} --- categorical and quantitative variables and their
displays, describing a distribution in context, mean, median, standard deviation, IQR and
resistance, the five-number summary and boxplots, comparing two distributions, and today's normal
model. Then \emph{Unit~2: Exploring Two-Variable Data} --- every question so far has been about
one variable at a time; next unit we put two measurements on the same individual and ask whether
they move together.
\end{skillbox}

% ── Teacher notes — one per component, in packet order (four: there is no activity) ──
\begin{teachernote}[Warm-Up]
Four minutes, five at the outside. Item~1 is a 45-second spiral --- do not let the boxplot
conversation reopen. Item~2 is mechanical and exists so that dividing by the standard deviation
is not the obstacle inside notes row~2; if the room is quick, take 60 seconds and move on.
\textbf{Item~3 is the one to debrief aloud}, and the sentence you want on the board in the
students' own words is \emph{``it is even on both sides, so split the leftover in half.''} Leave
it there all period --- notes problems 16 and~17, the homework, and the AP practice all reach
back for it. Item~3's percentages are 90 / 10 / 5 rather than the real 95 / 5 / 2.5 on purpose:
the real numbers arrive in notes problem~7 out of the students' own band counts, and they land
harder for having been computed than announced.
\end{teachernote}

\begin{teachernote}[Guided Notes \& Practice]
\textbf{This one document is the lesson --- pace it 24 / 10, then 8 for the debrief, then 8 for
the homework start.} The notes are a \emph{Main Ideas / Notes} table, not prose: five rows, 20
short problems, and six worked blocks you do at the board. \textbf{In every row you narrate the
prose, fill the lines, work the block, and work the first problem; the class works the rest with
the pen in their hand.} Row~1 is quick --- ten fills, three labelled answers --- and the only
thing worth slowing for is putting $\mu$, $\sigma$, $\bar x$, $s_x$ on the board in two columns
and leaving them there.

\textbf{Row~2 is where the arithmetic becomes an idea.} The dotplots are given: point at the two
gold dots and say out loud that 12.0 metres and 12.3 seconds are not on one scale. Work both
$z$-score blocks yourself so the subtraction and the division are settled in public, then release
4--6. \textbf{Problem~5 is the one to read while circulating} --- a student who writes ``1.5
below'' without saying that below is \emph{good} in a sprint has half the idea.

\textbf{Row~3's problem~7 is the table that earns the empirical rule --- do not announce
68--95--99.7 before it is filled.} The two records are deliberately unrelated: you work 341/476/499
of 500 loaves in the block, and they do 273/381/399 of 400 temperature readings, which come out
68.2, 95.2, 99.8 and 68.25, 95.25, 99.75. Write their six numbers on the board \emph{first}, then
round them into the rule in problem~9. Problem~8 is the sentence that matters: shape, not units
and not sample size, is what fixes the percentages.

\textbf{Row~4 is the crux, and you work none of it.} Read the coach's argument aloud exactly as
written and take a hand count on problem~11 before anyone speaks; the split in the room \emph{is}
the lesson. Problem~12 must end with a named athlete and a reason that mentions direction.
Problem~13 is the companion error --- negative read as worse --- and it is the one that survives
into the homework's item~5. Problem~14 is the counterweight and the AP-tested half of the rule.

\textbf{The Bread Loaf is where the pen changes hands.} Problems 15--20 on the bakery the notes
already pictured, worked entirely together, with its three blocks (834 g, the halved tail, 754 g)
worked \emph{with} the class. Problems 16 and~17 are the halve-the-tail move from warm-up~3 in its
real numbers; \textbf{problem~19 is the judgment problem and 20 the honest limit}, and they are
what to protect if the clock is tight --- 15, 16, 19, 20 is the right subset.

\textbf{There is no independent practice set on this page, and that is deliberate --- the
homework is the individual practice.} Guided Practice is the last thing on the page; the period
ends with students starting the homework in front of you, which is where your formative read
comes from. If more than a third of the room misses homework item~3, \textbf{reopen Maya and Dev
in the unit-test review}, and say so plainly.

\textbf{Debrief (8 min), spoken.} Open on the \emph{almost}-right problem~12 --- both $z$-scores
computed, the direction dropped --- because that is the exact gap the exam punishes. If time is
short, cut problem~14 to one sentence; making a student state a full $z$-score interpretation in
context is the move that cannot be cut.
\end{teachernote}

\begin{teachernote}[AP Practice]
\textbf{This page is not required and not scored} --- the graded homework is the section behind it
at the back of the packet. Work it as AP-format reps and go over it in the unit-test review.
\textbf{MC item~4 is the one worth reading closely}: its correct answer is that the empirical rule
\emph{does not apply}, because the length-of-stay histogram is strongly right-skewed. Students who
have just spent a period being handed 68--95--99.7 will apply it to anything with an axis, and the
exam knows that. If more than a handful miss it, put the notes' bread histogram back on the screen
beside the skewed one and ask what is missing --- that is notes problem~14 again, in AP wording.
On the free response, part (d) is the highest-value part --- it asks whether the model is
appropriate at all, which is Skill 2.D in the form the exam actually asks it, and it is the part
students most often skip. The birth-weight record appears nowhere else in the packet, so treat
what comes back as evidence that the notes transferred.
\end{teachernote}

\begin{teachernote}[Homework]
\textbf{This is the graded page --- scored out of 10, 2 points an item, and due the first class
after two study halls.} The eight minutes at the end of the period are a \emph{start}, not a
completion: put them on 1, 2, and~3 and let 4 and~5 go home. Item~3 is the one to read first when
scoring; ``Ben, he scored 88'' means the lesson did not land for that student, and it is worth two
minutes with them individually before the unit test. Item~2's third value is the diagnostic: a
student who writes ``1.5'' instead of ``$-1.5$'' has dropped the sign, and the sign is half the
idea. Item~5 is notes problem~13 in AP clothing --- the direction reversed, in a race where a
lower time is better, and distractor (C) reaches the right runner by the wrong reasoning; sort it
into three piles when you score (chose (B) with both $z$-scores; chose (B) with the raw times;
chose (A), (C), or (E)). Hold the context line while scoring --- an interpretation that says
``$z = -1.5$'' without naming heart rate, the mean of 74, and the direction is incomplete,
exactly as it would be on the exam.

\textbf{Because this is the last content lesson of the unit}, the closing preview box points at
the unit test and Unit~2; say the test date out loud when you assign the page, since the box alone
will not do it. \textbf{Paper is the default}; on the occasions you assign DeltaMath instead,
strike the score cell on the cover and tell students this page is now optional practice. Do not
assign both.
\end{teachernote}
\end{document}
